%%%%%%%%%%%%%%%%%%%%%%%%%%%%%%%%%%%%%%%%%%%%%%%%%%%%%%%%%%%%%%%%%%%%%%%%%%%%%%%
%                                                                             %
%           TEMPLATE LATEX PER TESI                                           %
%           ______________                                                    %
%                                                                             %
%           Ultima revisione: 10 Dicembre 2025                                %
%           Revisori: G.Presti; L.A.Ludovico; F. Avanzini; M. Tiraboschi      %
%                                                                             %
%%%%%%%%%%%%%%%%%%%%%%%%%%%%%%%%%%%%%%%%%%%%%%%%%%%%%%%%%%%%%%%%%%%%%%%%%%%%%%%

% Usare l'opzione twoside se volete stampare fronte-retro. Es.:
% \documentclass[twoside,12pt]{report}
% Si consiglia di diminuire il font per tesi di dottorato in formato 17x24. Es:
% \documentclass[11pt]{report}
\documentclass[12pt]{report}

% --- PREAMBOLO ---------------------------------------------------------------
% Inserire qui eventuali package da includere o
% definizioni di comandi personalizzati...

% ...ad esempio il seguente package permette di inserire note visibili nel pdf di output,
% (il comando che lo precede serve a lasciare spazio alle note, va rimosso prima della consegna)
% \setlength {\marginparwidth }{2cm}
% \usepackage{todonotes}  % istruzioni sull'uso in sezione 3.2.9

% Selezione lingua
\usepackage[italian]{babel}

\usepackage{tesi}
% Puoi usare il font di default di LaTeX con la relativa opzione del package
% \usepackage[defaultfont]{tesi}
% Esiste anche un'opzione per il formato 17x24 per le tesi di dottorato
% \usepackage[phd]{tesi}

% In caso il copia-incolla del PDF generato perda gli spazi,
% provare a decommentare la seguente riga
% \pdfinterwordspaceon

% !!! INFORMAZIONI SULLA TESI DA COMPILARE !!!
\usepackage{pifont}
\usepackage{xcolor}
\newcommand{\cmark}{\textcolor{green!60!black}{\ding{51}}}
\newcommand{\xmark}{\textcolor{red!70!black}{\ding{55}}}
%   UNIVERSITA' E CORSO DI LAUREA:
\university{Università degli Studi di Milano}
\unilogo{immagini/loghi/unimi}
\faculty{Facoltà di Scienze e Tecnologie}
\department{Dipartimento di Informatica\\Giovanni Degli Antoni}
\cdl{Corso di Laurea Triennale in\\Corso di Laurea}

%   TITOLO TESI:
\title{A multimodal knowledge distillation model for emotive EEG characterization}

%   AUTORE:
\author{Jacopo Fichera}
\matricola{123456} % TODO
% "Elaborato Finale" per i CdL triennali
% "Tesi di Laurea" per i CdL magistrali
\typeofthesis{Elaborato Finale}

%   RELATORE E CORRELATORE:
\relatore{Prof.ssa Raffaella Lanzarotti}
\correlatore{Dot. Jacopo Burger}

%   LABORATORIO:
% Questa sezione crea una pagina di chiusura della tesi con
% il logo dell'ente/laboratorio presso cui si è svolto il tirocinio.
% Più afferenze/url/loghi sono supportate,
% e la frase può essere personalizzata.
% Qui trovate alcuni predefiniti del nostro dipartimento
% \adaptlab
% \aislab
% \anacletolab
% \bisplab
% \connetslab
% \everywarelab
% \falselab
% \iebilab
% \islab
% \lailalab
% \lalalab
% \lawlab
% \laserlab
% \limlab
% \mipslab
% \optlab
% \phuselab
% \ponglab
% \sesarlab
% \spdplab

% Esempio di personalizzazione della pagina di chiusura
% (non consegnate con questo esempio!)
% (da commentare in caso sia sufficiente una delle macro precedenti)
\lab{Laboratorio di Ricerca}
\lab[in collaborazione con l']{Azienda Specifica}
\laburl{https://di.unimi.it/it/ricerca/risorse-e-luoghi-della-ricerca/laboratori-di-ricerca}
\lablogo{immagini/redqmark}

% Con questo comando si può cambiare la dimensione (massima
% altezza e larghezza consentite) dei loghi
% \setlength\lablogosize{25mm}

%   ANNO ACCADEMICO
% \the\year inserisce l'anno corrente
% per specificare manualmente un anno accademico
% NON inserire nel formato 1970-1971, ma
% inserire solo 1970
\academicyear{\the\year}

%   INDICI:
% elenco delle figure (facoltativo)
% \figurespagetrue
% elenco delle tabelle (facoltativo)
% \tablespagetrue
% prefazioni nell'indice (facoltativo)
\prefaceintoctrue
% indice nell'indice (facoltativo)
\tocintoctrue
% --- FINE PREAMBOLO ----------------------------------------------------------

\begin{document}
% Creazione automatica della copertina
% Centra la8 copertina nel foglio: usa questo comando per la copertina esterna
% \makecenteredfrontpage
% Copertina allineata alle altre pagine: usa questo comando per la copertina interna
    \makefrontpage

%
%			PAGINA DI DEDICA E/O CITAZIONE
%			facoltativa, questa è l'unica cosa che dovete formattare a mano, un po' come vi pare
%

    {\raggedleft \large \sl Questo lavoro \`{e} dedicato ai miei genitori\\

    \vspace{2cm}

    ``What I cannot create, I do not understand''

    \bigskip

    \--- Richard Feynman\\

    \vspace{2cm}

    ``It's not only powerful,\\but it's also inadequate''

    \bigskip

    \--- Miller Puckette\\}

    \clearpage
    \beforepreface

%
%			PREFAZIONE (facoltativa)
%

% \prefacesection{Prefazione}
% Le prefazioni non sono molto comuni, tuttavia a volte capita che qualcuno voglia dire qualcosa che esuli dal lavoro in sé (come un meta-commento sull'elaborato), o voglia fornire informazioni riguardanti l'eventuale progetto entro cui la tesi si colloca (in questo caso è probabile che sia il relatore a scrivere questa parte).

%
%			RINGRAZIAMENTI (facoltativi)
%

    \prefacesection{Ringraziamenti}
    Questa sezione, facoltativa, contiene i ringraziamenti.

%
%			Creazione automatica dell'indice
%

    \afterpreface
    %! Author = jacopo
%! Date = 3/12/26

% todo add citations
\chapter{Introduction}
\label{cap:introduzione}
Machine Learning advancements can be applied across a wide range of domains due to their ability to model complex data.
This has also influenced the field of affective computing, where tasks often rely on electroencephalography (EEG) to infer emotional and cognitive states.
These signals  are among the most commonly used physiological measurements in the domain, as they provide one of the most direct proxies currently available for capturing such states. \cite{RAHMAN2021104696}

In recent years, many models have been developed for EEG representation learning.
These general-purpose models, \textit{foundation models} \cite{Schneider2024}, in the field tend to be limited by various reasons, with the low cardinality of EEG-centric datasets being one of them.
This is especially noticeable once one considers the size of corpora of most text-based state-of-the-art (SotA) models.
The reason for this is that textual corpora can be collected at large scale using publicly available resources such as the internet.
In contrast, EEG data gathering has more requirements.
Most experimental setups require researchers and specific equipment.
They also require significant time both from researchers and participants.
Many datasets also do not follow common acquisition standards such as the 10–20 electrode placement system making them at times hard to compare or interchange in tasks.

Although multimodal foundation models have shown success in vision and language, their application to EEG representation learning remains limited.
Among EEG models in the affective field, only a small number of models manage to integrate additional modalities.
\textit{BIOT} \cite{NEURIPS2023_f6b30f3e} is one of those, but its approach is still constrained to the physiological domain.
Integration of modalities such as video and audio is even less common in EEG-centered models, with \textit{Hyper-MML} \cite{kang2025hypergraphmultimodallearningeegbased} being one of the few examples.
These modalities have the advantage of receiving greater attention in the broader machine learning community as they are used for a wider range of applications.
Consequence is that more discoveries and better-performing systems are constantly being made in these contexts.
\FA[ok]{Qui puoi entrare un po' più nel dettaglio rispetto al gap presente in letteratura: i lavori esistenti non sfruttano la possibilità di trasferire informazione, in setting affettivi, da modalità dense di dati a modalità target più scarse - gli EEG.}
With this problem in mind, the goal of the thesis is to leverage cross-modal knowledge to integrate information into EEG representations. \FA[ok]{e qui di conseguenza entra un po' più nel dettaglio su quale sia l'obiettivo della tesi, ovvero ottenere un backbone forte multimodale affective bla bla}
This is achieved by a strong backbone for affective multi-modal processing that overcomes limited data and takes advantage of stronger representational models from other, richer, modality contexts.
This is based on the assumption that the same phenomena are encoded in different but complementary modalities \cite{AHMED2023200171}.
For example, the emotional states reflected in the EEG signals might be correlated with facial expression, speech characteristics, and/or body movements.
Leveraging such cross-modal relationships may help improve EEG representations when EEG data alone are limited.

The proposed \textit{EEGAVI} transfers information from multiple modalities into the EEG, aiming to enhance the representational capacity.
It employs a principled integration of existing techniques (cross-attention + contrastive loss + knowledge distillation) applied to EEG under data scarcity.
It takes inspiration from \textit{Flamingo} \cite{alayrac2022flamingovisuallanguagemodel} that leverages video-textual relations by using frozen encoders.
EEGAVI does not interleave the main modality's frozen layers with cross-attention like google's model but it rather
makes these steps sequential and builds on top of the encoders a multi-modal cross-attention mechanism.
For the EEG data stream, \textit{CBraMod} \cite{wang2025cbramodcrisscrossbrainfoundation} is the selected backbone.
The modality acts as pivot in the architecture being always required.
Other modalities are instead considered supports and grant additional information in the forward pass of the model.
This way the model fuses the information into the pre-computed EEG embeddings.
It is achived by adopting a contrastive learning paradigm on different affective datasets in pretraining: \textit{EAV} \cite{article-eav}, \textit{AMIGOS} \cite{8554112}, \textit{DEAP} \cite{5871728} and \textit{MAHNOB-HCI} \cite{5975141}. % todo citations

The model is so able to learn both to leverage additional input streams in runtime and to restructure the input sequence in order to make it more expressive based on pretraining correlations of the modalities.
\FA[ok]{qui sarebbe meglio mettere in evidenza e separare la pipeline, aggiungendo anche un po' più nel dettaglio il tipo di dati usati e che le EEG sono la modalità pivot, da quelli che sono i contributi originali della tesi, magari partendo da flamingo aggiungendo la pipeline nostra e poi sottolineando i contributi, magari come elenco puntato (se ti piace, non è obbligatorio)}

Beyond proposing a multi-modal foundation model, the research question focuses on the impact of \textit{knowledge distillation} in the proposed context.
In particular, the goal of the investigation is to determine whether distillation in a context where sources and target do not belong to the same modality is of benefit to the training pipeline or not.
To answer this question, an instance of \textit{EEGAVI} is trained and its performance evaluated against CBraMod as a reference baseline.
\FA[ok]{qui si può spiegare un po' meglio come abbiamo fatto il test}
The tests are performed as linear probes \cite{alain2018understandingintermediatelayersusing}.
A test is conducted on the \textit{FACED} \cite{Chen2023} dataset to assess the integrity of EEG embeddings when there are no supporting modalities.
Other linear probes are performed intra- and cross dataset to assess the gains of multi-modal fusion.
% todo questa parte da finire quando so esattamente cosa dire
Lastly, the model is trained on a specific task on the \textit{VR} dataset to see a possible real world application of the model. % TODO Citazioni che mancano


The proposed approach is limited to the affective aspect of neurological data processing, but the core idea can be applied to any field. \FA[ok]{Questo meglio metterlo alla fine come conclusione}

    %! Author = jacopo
%! Date = 3/12/26

% Preamble
\documentclass[11pt]{article}

% Packages
\usepackage{amsmath}

% Document
\begin{document}



\end{document}
    %! Author = jacopo
%! Date = 3/12/26


\section{EEGAVI}
\label{chap:eegavi}

\textit{EEGAVI} is a multimodal architecture built on frozen backbone encoders for each modality.
The frozen representations are further elaborated before getting to a \textit{cross-attention} mechanism.
\textit{EEG} is the main modality of the model acting as a \textit{pivot} while the others are \textit{supports}.
This structure allows the model to generalize well even with limited or absent supporting modalities.
Like many other representation models, \textit{EEGAVI} outputs a \textit{CLS} token that condenses sample information in a tighter space.
It is used during training for contrastive alignment.
The model's default implementation is able to process EEG, video, text, audio and ECG data.

% todo expand further

\subsection{Data}
For the development of the model, we use data from different multimodal datasets.
Most of these lack some modalities that the model can process.

\begin{table}
    \centering
    \begin{tabular}{c | c | c| c | c| c| c | c| c}
        Dataset & Subjects & Experiments & EEG    & Video  & Audio  & Text   & ECG    & Samples \\
        \hline
        AMIGOS  & 40       & 16 + 4      & \cmark & \cmark & \cmark & \cmark & \cmark & 22k     \\
        EAV     & 42       & 200         & \cmark & \cmark & \cmark & \cmark & \xmark & 40k     \\
        DEAP    & 32       & 40          & \cmark & \cmark & \xmark & \xmark & \xmark & 12k     \\
        MAHNOB  & 29       & 20(1*)      & \cmark & \cmark & \xmark & \xmark & \cmark & 1k      \\
    \end{tabular}
    \caption{
        Datasets used for training summary.
        MAHNOB is missing most experiments as only few for participant are accessible.
    }
    \label{tab:datasets}
\end{table}


\begin{itemize}
    \item \textbf{AMIGOS}: The dataset was composed for the research on affect, personality traits and mood by means of neuro-physiological signals.
    The experiments were performed on a total of 40 participants and are divided in two different subsets.
    For the first part, 16 short emotional videos from famous movies were shown to each participant alone.
    For the second one, subjects watched 4 long videos, some alone and some in groups.
    Self-assessments and external annotations on affective levels and emotions were also tracked.

    \item \textbf{EAV}\cite{article-eav}: \textit{EEG-Audio-Video Dataset for Emotion Recognition in Conversational Contexts (EAV)}
    measures 30-channel EEG, audio and video from 42 participants.
    The scenarios proposed are done in a cue-based conversation designed around five distinct emotions.
    Each participant contributed in a total of 200 interactions.
    Participants tracked self-assessments for the emotion of the experiment (happy, calm, angry and sad) and their level of arousal.
    They were also evaluated by the experimenters on the same protocol.

    \item \textbf{MAHNOB-HCI}\cite{5975141}: is a broad dataset recorded in response to affective stimuli.
    Unlike other datasets \textit{MAHNOB} also tracks gaze signals, but it is too limited in size to use.
    We did not include an explicit gaze stream due to a lack of calibrated gaze targets and potential estimation noise from frontal videos.
    However, eye- and head-related cues could be extracted as an auxiliary modality in future work.

    \item \textbf{DEAP}: \textit{Database for Emotion Analysis} is a dataset composed of data from 32 participants.
    Each watched 40 one-minute long clips of music videos.
    Subjects rated their valence, arousal, like/dislike, dominance and familiarity levels.
    Audio for these is not given for copyright reasons.

\end{itemize}


\begin{table}
    \centering
    \begin{tabular}{c | c | c | c | c}
        Dataset & Train (70\%) & Validation (10\%) & Test (15\%) \\
        \hline
        AMIGOS  & 30           & 4                 & 6           \\
        EAV     & 32           & 4                 & 6           \\
        DEAP    & 24           & 3                 & 5           \\
        MAHNOB  & 22           & 3                 & 4           \\
    \end{tabular}
    \caption{
        Partitioning of the subjects in the three data splits for training \textit{EEGAVI}.
    }
    \label{tab:datasetspartitoning}
\end{table}

All datasets are partitioned in train, validation and test set around (70-10-15)\%.
As commonly done for EEG applications, the partitioning is not performed on the sample level but on the participants directly,
meaning the real split proportions might be a little off if data is corrupted, missing or a participant just made fewer experiments.
This should allow the model to force a stronger cross subject generalization.
Partitioning of the data is reported in table \ref{tab:datasetspartitoning}.
%todo rework with new infos

\subsubsection{Sampling Strategy}
%todo rework with new infos
Affective responses unfold across multiple temporal scales, for this reason the model must generalize well on both short events and longer temporal dynamics.
While short clips may capture localized affective outbursts, long ones can encode the evolution of a subject's perception over time.
For the model to be robust across scales we sample variable-length EEG segments between 1 and 32 seconds using an EEG-centric sampling strategy.

Segment durations are sampled from a log-uniform distribution, favoring shorter clips.
A sample is either taken at random from the clip with probability $p$ or selected as a point of interest from the signal with probability $1-p$.
In the experimental setup $p=0.6$.
When an \textit{anchor} point is selected short segments are built by left-aligning the window at the anchor time rather than centering it.
This design may introduce a slight potential positional bias in short clips; however, longer segments are centered on previously extracted short ones and therefore do not have this behavior.

If the resulting segment respects the overlap constraints defined for that class it is stored as a sample, else discarded.
Segments are of three different settings as reported in table \ref{tab:sample}.
The process is iterated until it fails to present new valid segments after running out of patience.

\begin{table}
    \centering
    \begin{tabular}{c|c|c|c|c}
        %todo update
        Class  & Max length (s) & Same class IoU & Cross class IoU & Jitter Frac \\ \hline
        Short  & 4              & 0              & 0.7             & -           \\
        Medium & 16             & 0.35           & 0.7             & 0.3         \\
        Large  & 32             & 0.25           & 0.7             & 0.4         \\
    \end{tabular}
    \caption{The jitter fraction indicates the maximum relative displacement applied to a segment's center. }
    \label{tab:sample}
\end{table}

The candidate extractor first computes STFT-based frame features per EEG channel.
The points of interest are extracted in the STFT domain following a set of heuristics commonly used when processing EEG.
These anchors are identified through a multi-descriptor scoring mechanism that combines the following indicators:
\begin{itemize}
    \item \textit{Relative band-power per frame}: Fraction of total spectral power contained in predefined frequency bands (0.5–4, 4–8, 8–13, 13–30 Hz).
    Sudden dominance of a specific band indicates a change in oscillatory regime contributing positively to salience.

    \item \textit{Spectral entropy}: Shannon entropy of the normalized power spectrum; high when the spectrum is flat (distributed energy), low when concentrated.
    Lower entropy is favored as it indicates structured spectral content rather than broadband noise.

    \item \textit{Spectral flux}: Frame-to-frame positive spectral change computed on magnitude spectra, capturing rapid spectral dynamics.
    Rapid spectral changes are rewarded in the saliency extraction process.

    \item \textit{Root-mean-square (RMS)}: Measures the average signal energy in a window, independent of frequency content.
    Higher energy contributes positively to candidate selection as it highlights strong neural responses.

    \item \textit{Line length}: It is a simple but powerful time-domain EEG feature that measures the total amount of change in the signal within a window.
    It computes the average absolute first difference, acting as a proxy for signal complexity.
    In EEG it is often associated with high activity as seizures/epilepsy or motion artifacts, therefore its contribution is moderated in the scoring process.

    \item \textit{Teager–Kaiser Energy Operator (TKEO)}: It captures bursty, high-frequency energy.
    Sensitive to rapid changes, it is popular in EEG analysis; adds robustness to the process as it might capture patterns not recognized by RMS alone.
\end{itemize}

Afterward, each channel's features are standardized with a rolling median/MAD estimator with a 30s window.
This is done to prevent high-variance features from dominating the saliency score, mitigating the effect of outliers.
A saliency score per channel as a weighted sum of z-scored features is formed.
These scores are aggregated by taking the mean of the top $k$ channels.
Artifact annotations are used to downweight candidate frames before anchor selection.
Lastly, a greedy top-score selection with a minimum time gap constraint is applied to select the points.

    %! Author = jacopo
%! Date = 3/12/26


\chapter{Test}
\label{cap:test}


\section{Hyperparameter Tuning}

\subsubsection{Exploratory Search}

\subsubsection{Refinement}

\subsubsection{Final Selection}
full trained on data
    %! Author = jacopo
%! Date = 3/12/26

% Preamble
\chapter{Conclusions}
\label{ch:conclusions}



    \section{Sviluppi futuri}

    Tra gli sviluppi futuri in genere trovano posto quelle migliorie non realizzate per mancanza di tempo, la cui necessità è emersa solo dopo i test, e che riguardano il progetto ad un livello di astrazione più alto (nel caso di tesi che si inquadrano in una linea di ricerca).


%
%			APENDICE: materiali aggiuntivi e dimostrazioni
%

    \appendix


    \chapter{Documenti da Produrre}
    I membri del LIM in generale non richiedono agli studenti di consegnare una copia cartacea del proprio lavoro. In altri casi, si invita lo studente a concordare questo punto con il proprio relatore.
    In prossimità della discussione viene richiesto loro di condividere in formato digitale:
    \begin{itemize}
        \item la versione finale dell'elaborato;
        \item il riassunto caricato in piattaforma;
        \item il codice sviluppato;
        \item tutti i materiali di supporto utilizzati e/o prodotti;
        \item la presentazione predisposta per la discussione.
    \end{itemize}

    La consegna può avvenire via chiavetta USB o attraverso i sistemi di condivisione dei file (WeTransfer, Dropbox, Google Drive, ecc.).


    \section{Riassunto}
    In prossimità della data di laurea, viene richiesto agli studenti di caricare in piattaforma anche il riassunto. La sua funzione è permettere alla commissione di visionare rapidamente gli argomenti trattati da ciascun elaborato.

    L'estensione di tale documento non dovrebbe superare le due pagine, e spesso una pagina è più che adeguata allo scopo. Riguardo i contenuti, il riassunto riprende -- in maniera molto succinta -- lo schema dell'elaborato finale, e quindi, ad esempio, la scaletta proposta in \ref{sec:organizzazione}.

    Visto lo scopo del documento, è buona prassi bilanciare il testo in modo da dedicare poche righe allo stato dell'arte, in favore delle parti originali dell'elaborato.


    \section{Presentazione}
    Gli studenti delle lauree triennali hanno a disposizione al massimo 10 minuti per presentare il proprio elaborato, che diventano 15 per le lauree magistrali.

    Si raccomanda di provare ripetutamente il discorso cronometrandosi e considerando i tempi morti per l'accensione del PC, l'avvio dei software, il passaggio da un'applicazione a un'altra, e via dicendo. Si segnala inoltre la possibilità di prenotare l'aula della discussione per provare in anticipo i collegamenti e verificare che immagini e audio vengano correttamente riprodotti.

    Per via degli stringenti vincoli temporali, difficilmente una presentazione multimediale può superare la soglia di 15 slide. Questo valore è puramente indicativo, perché dipende dalla densità di informazioni della slide e dal tempo dedicato a ciascuna slide mentre si sta presentando.

    La scelta tecnologica più comune è Microsoft PowerPoint, ma qualsiasi alternativa (Prezi, Acrobat PDF, \LaTeX, ecc.) è ammissibile.

    Gli errori più comuni da evitare sono:
    \begin{itemize}
        \item slide dal contenuto testuale troppo ricco;
        \item slide che verranno lette parola per parola in fase di presentazione.
    \end{itemize}

    Entrambi i problemi si risolvono impostando il testo in maniera estremamente schematica, tendenzialmente per elenchi puntati. Si presti attenzione, però, a usare una forma omogenea.
    Ad esempio, sarebbe poco elegante il seguente elenco, che include aggettivi, sostantivi e frasi di senso compiuto, e non fa un uso corretto di iniziali maiuscole:
    \begin{itemize}
        \item applicativo esteticamente molto gradevole;
        \item C'è la possibilità di interagire con altri utenti;
        \item l'utente può esportare file in formato XML.
    \end{itemize}

    Una possibile revisione, basata solo sui sostantivi, è:
    \begin{itemize}
        \item qualità estetica;
        \item interazione tra utenti;
        \item supporto dell'XML.
    \end{itemize}

    Un altro consiglio per rendere la presentazione cognitivamente più gestibile da parte del pubblico è quello di aggiungere in calce ad ogni slide un indicatore dello stato di avanzamento (ad esempio ``Slide 3 di 15''). In questo modo, in caso di esaurimento del tempo concesso la commissione saprà se dare modo di concludere (``Slide 14 di 15'') o se fermare il candidato (``Slide 14 di 138''). In assenza di un indicatore la commissione tende a presupporre il secondo caso.

    Per quanto riguarda gli argomenti da trattare durante la discussione, è opportuno esporre brevemente lo stato dell'arte per concentrarsi sul proprio lavoro, sui test e sui risultati. Una tecnica per capire quanto tempo dedicare a ogni argomento consiste nel dare un peso a ciascuna parte di cui si vuole parlare, quindi assegnare un numero di minuti proporzionale a detto peso in modo da raggiungere il tempo massimo a disposizione.

    Video dimostrativi e \textit{live demo} sono gradite, ma è sconsigliabile dedicare troppo tempo a questi contributi, a meno che al contenuto multimediale si sovrapponga una spiegazione da parte del candidato.


%
%			BIBLIOGRAFIA
%

% Si può specificare a che livello della TOC deve essere la bibliografia.
% Il default è 'chapter', per 'part' usare
% \beforebibliography[part]
    \beforebibliography
    \bibliographystyle{unsrt}
    \bibliography{bibliografia}

% Pagina di chiusura
    \closingpage

\end{document}
